% ====================================================================
\section{Introduction}\label{sec:intro}
% ====================================================================

\subsection{The problem and the results}

Let $T=(V,E)$ be a finite simple undirected tree in which every vertex carries a bit. In
\emph{synchronous majority dynamics} all vertices update at the same time: a vertex takes the strict
majority of the bits of its neighbours, and keeps its own bit when the neighbours are tied. There are
no weights, no directed edges and no physical self-loops. We study two decision problems for this
dynamics on trees of maximum degree four (formal definitions are in \cref{sec:prelim}):
\begin{itemize}
\item $\Pred$: given the tree, an initial configuration, a vertex $v^\ast$ and a time $\tau$ written
  in binary, is $v^\ast$ in state $1$ at time $\tau$?
\item $\Ever$: given the tree, an initial configuration and a vertex $z$ initially in state $0$, is
  $z$ in state $1$ at some time?
\end{itemize}

The heart of the paper is the following simulation theorem, stated precisely as
\cref{thm:simulation}. Monotone Boolean straight-line programs are defined in \cref{sec:model}.

\begin{mainthm}[Simulation theorem]\label{thmA}
There is a logspace transducer that maps every monotone Boolean straight-line program $\mathcal P$
to a tree $T$ of maximum degree four, an initial configuration, a vertex $v^\ast$ and a time $\tau$
written in binary, such that the state of $v^\ast$ at time $\tau$ under synchronous majority equals
the final value of the designated register of $\mathcal P$. If $\mathcal P$ has $m$ registers in
single-assignment normal form, then $\tau=O(m^2)$, $|V(T)|=O(m^{13})$, and all vertex degrees lie in
$\{1,3,4\}$. A variant of the construction produces, instead of $(v^\ast,\tau)$, a vertex $z$
initially in state $0$ which is ever in state $1$ if and only if the program outputs $1$.
\end{mainthm}

Combined with the classical $\Ptime$-completeness of monotone circuit value (\cref{sec:pcomplete})
and with the classical bounded-transient behaviour of symmetric majority networks
(\cref{sec:membership}), the simulation theorem gives the two completeness results.

\begin{mainthm}\label{thmB}
$\Pred$ is $\Ptime$-complete under logspace many-one reductions. Hardness already holds for trees
whose vertex degrees all belong to $\{1,3,4\}$, with a query time polynomial in the instance size.
\end{mainthm}

\begin{mainthm}\label{thmC}
$\Ever$ is $\Ptime$-complete under logspace many-one reductions, already on trees whose vertex
degrees all belong to $\{1,3,4\}$.
\end{mainthm}

Since trees have treewidth one, hardness holds at treewidth one. The results say nothing about
maximum degree three; see \cref{rem:scope}.

\subsection{Architecture of the proof}\label{sec:architecture}

The proof has six parts, each in its own section.
\begin{enumerate}[label=(\arabic*)]
\item \emph{Computational model} (\cref{sec:model}). Monotone Boolean straight-line programs:
  Boolean registers, $\AND$/$\OR$ instructions, unrestricted reuse of stored values, and no
  conditional jumps. A logspace normal form gives single assignment and topological order.
\item \emph{$\Ptime$-completeness of program evaluation} (\cref{sec:pcomplete}). This is
  essentially the classical result of Ladner and Goldschlager. We do not reprove it.
\item \emph{Simulation theorem} (\cref{sec:simulation}). We state it and split its proof into
  two layers. The first is a \emph{controlled processor}, whose control vertices follow prescribed
  words. The second consists of \emph{autonomous clocks} that generate those words. A substitution
  principle (\cref{prop:substitution}) glues the layers.
\item \emph{Modular simulation lemmas}:
  \begin{itemize}
  \item storage, transport, logic and reading, as local contracts (\cref{sec:modules});
  \item clocks and synchronisation, through a phase calculus and robust generators
    (\cref{sec:clocks});
  \item restoration and composition, through explicit invariants and induction over the program
    (\cref{sec:processor}).
  \end{itemize}
\item \emph{Reduction analysis} (\cref{sec:analysis}): correctness, size, time, degree, acyclicity
  and logspace constructibility.
\item \emph{Completeness results} (\cref{sec:completeness}). Membership in $\Ptime$ is proved
  separately, and eventual activation is treated as an additional result with its own
  finite-window reader.
\end{enumerate}

\Cref{fig:dependencies} shows how the results depend on each other. An arrow $X\to Y$ means that the
proof of $Y$ uses $X$.

\begin{figure}[ht]
\centering
\begin{tikzpicture}[
  box/.style={draw,rounded corners=2pt,align=center,font=\scriptsize,inner sep=2.5pt,
              text width=2.75cm,minimum height=9mm},
  cls/.style={box,fill=gray!14},
  main/.style={box,fill=blue!6,text width=2.2cm,minimum height=6mm},
  arr/.style={-{Stealth[length=1.8mm]},semithick,gray!65!black}]
\def\xa{0}\def\xb{3.7}\def\xc{7.4}\def\xd{11.1}
% row 1: local modules and clocks
\node[box] (cell)  at (\xa,0) {cells, junctions,\\ duality, frame\\ \cref{lem:modes,lem:duality,lem:frame}};
\node[box] (prims) at (\xb,0) {\prim{SET}, \prim{COPY}, \prim{MOVE},\\ \prim{ROUTE}, \prim{GATE}, \prim{READ}\\ \cref{sec:modules}};
\node[box] (phase) at (\xc,0) {phase calculus\\ \cref{lem:cellphase,lem:concat,lem:timing}};
\node[box] (gens)  at (\xd,0) {robust generators\\ \cref{thm:phasedefect,cor:words}};
% row 2: composition
\node[box] (proc)  at (\xb,-1.9) {milestones, instruction\\ step, normalisation\\ \cref{prop:instruction,lem:normalize}};
\node[box] (forced) at (\xc,-1.9) {forced-control\\ simulation\\ \cref{thm:forced}};
\node[box] (subst) at (\xd,-1.9) {substitution\\ principle\\ \cref{prop:substitution}};
% row 3: the simulation theorem
\node[box] (nf)    at (\xa,-3.8) {normal form\\ \cref{lem:normalform}};
\node[box] (logsp) at (\xb,-3.8) {logspace transducer\\ \cref{prop:logspace}};
\node[box,very thick] (sim) at (\xc,-3.8) {\textbf{simulation theorem}\\ \cref{thm:simulation}};
\node[box] (reader) at (\xd,-3.8) {finite-window reader\\ \cref{lem:reader,lem:finalwrite}};
% row 4: classical inputs and main results
\node[cls] (mcvp)   at (\xa,-5.7) {ordered MCVP is\\ $\Ptime$-complete (classical)\\ \cref{thm:mcvp}};
\node[main] (B)     at (\xb,-5.7) {\textbf{\cref{thmB}}\\ prediction};
\node[cls] (energy) at (\xc,-5.7) {energy, period $\le2$\\ (classical)\\ \cref{prop:membership}};
\node[main] (C)     at (\xd,-5.7) {\textbf{\cref{thmC}}\\ activation};
% arrows
\draw[arr] (cell) -- (prims);
\draw[arr] (prims) -- (proc);
\draw[arr] (cell.south) |- (proc.west);
\draw[arr] (phase) -- (gens);
\draw[arr] (phase) -- (forced);
\draw[arr] (proc) -- (forced);
\draw[arr] (gens) -- (subst);
\draw[arr] (forced) -- (sim);
\draw[arr] (subst) -- (sim);
\draw[arr] (nf) -- (logsp);
\draw[arr] (logsp) -- (sim);
\draw[arr] (gens.east) -- ++(0.3,0) |- (reader.east);
\draw[arr] (sim) -- (B);
\draw[arr] (sim) -- (C);
\draw[arr] (reader) -- (C);
\draw[arr] (mcvp) -- (B);
\draw[arr] (energy) -- (B);
\draw[arr] (energy) -- (C);
\end{tikzpicture}
\caption{Dependencies between the main results. Shaded boxes are classical results that we cite.}
\label{fig:dependencies}
\end{figure}

\subsection{Classical results and original contributions}\label{sec:credits}

We cite the following results and do not reprove them, except where a short proof is needed in the
exact form we use.
\begin{itemize}
\item $\Ptime$-completeness of the circuit value problem for circuits presented as straight-line
  sequences \cite{Lad75}, and of its monotone restriction \cite{Gol77}; see also
  \cite[Appendix~A.1]{GHR95}. This is \cref{thm:mcvp}.
\item Closure of logspace reductions under composition \cite{Jon75,LL76}.
\item Period at most two and a linear transient for symmetric threshold networks, proved with a
  bounded energy function \cite{GO80,PS83,GFP85}. We include the short proof in
  \cref{prop:membership}, because we need the explicit transient bound for binary-encoded times.
\end{itemize}

The following constructions and results come from the original manuscript~\cite{Src}. To our
knowledge they are not in the classical literature:
\begin{itemize}
\item robust tree generators of all phase-defect words, with degree at most four and polynomial
  size, constructible in logspace (\cref{sec:generators});
\item the two-vertex \textsf{EDGE} memory cell with modes \RESET/\HOLD/\LOAD, and the local
  operations \prim{MOVE}, \prim{ROUTE} and the five-transition gates (\cref{sec:modules});
\item the fourteen-transition non-destructive reader that restores its isolating cells
  (\cref{lem:read});
\item the serial tree processor, its normal configurations, and the background normalisation wave
  (\cref{sec:processor});
\item the finite-window reader and the special final write used for eventual activation
  (\cref{sec:ever}).
\end{itemize}

This version reorganises the proof and adds the following organisational lemmas, each proved here:
\begin{itemize}
\item a mode-level duality lemma (\cref{lem:duality});
\item a frame lemma (\cref{lem:frame});
\item a cell-phase lemma (\cref{lem:cellphase});
\item a concatenation lemma that excludes the merging of control blocks (\cref{lem:concat});
\item a table of timing signatures of all primitives (\cref{tab:timing});
\item the substitution principle (\cref{prop:substitution});
\item an explicit list of milestone configurations (\cref{tab:milestones});
\item a closed formula for the observation time (\cref{prop:instruction,cor:tau}).
\end{itemize}
\Cref{app:map} records the correspondence with the original manuscript, and \cref{app:audit} lists
the points that the reorganisation made explicit.

\subsection{Conventions}\label{sec:conventions}

We write $\oplus$ for exclusive or. For an integer $t$, its \emph{parity} is $t\bmod 2$. For a
statement $S$, $\one\{S\}\in\{0,1\}$ is its indicator. For integers $a\le b$, $[a,b)$ is the set of
integers $t$ with $a\le t<b$.

A configuration at \emph{state time} $t$ is the vector of all bits at time $t$. A bit that a
vertex reads at state time $t$ is called a \emph{control value at control time $t$}: it influences
the transition from state time $t$ to state time $t+1$. In schedules, \emph{``during $[s,s+\ell)$''}
refers to control times. An operation occupying the control times $[s,s+\ell)$ takes the
configuration at state time $s$ to the configuration at state time $s+\ell$, and has
\emph{duration} $\ell$.
